\section{What additional information can language actually supply?}
\label{exec14:l4-development}

The preceding comparisons leave a practical design question. If an assistant
can explain a sensible causal story but the data identify only an equivalence
class, which additional information is the explanation contributing? This can
be answered precisely before choosing a model provider. A restriction extracted
from a dated policy document may eliminate a graph because the policy fixes the
order of decisions. A recollection about lending in general may merely express
a prior preference. A paraphrase of the same empirical association contributes
no new identifying information. The interface should make these three roles
visible to the analyst.

Write $P$ for the observed-data law and $\mathcal G(P;\mathcal A)$ for the set
of graphs compatible with that law under assumptions $\mathcal A$. A numerical
discovery procedure estimates features of this set. A text-derived restriction
$R$ instead proposes the smaller set
\[
 \mathcal G_R(P;\mathcal A)=\{G\in\mathcal G(P;\mathcal A):R(G)=1\}.
\]
This is a useful operation if the restriction is applicable and justified. It
is also a way to manufacture certainty when an erroneous restriction deletes
the true graph. The source, date, population and interpretation of $R$ are
therefore part of the scientific object. A statement that an account review
occurs before an offer concerns a time order; it does not by itself assert that
review has no other causes or that all such causes were measured.

For a deposit-pricing question, a policy may establish that a national rate is
announced before branch offers are made. That justifies a temporal restriction
for the relevant dates. It does not establish that managers did not anticipate
the announcement, that the national rate has no other effect on deposits, or
that branch offers were assigned independently of unrecorded local risk. A
language assistant should return the restriction actually supported by the
passage and list the stronger restrictions that would be needed for the desired
intervention calculation. This follows the distinction between an assumption
and its observable implications in graphical model construction
\citep{exec13_census_credible}.

An empty $\mathcal G_R$ is an informative conflict. It can result from an
incorrect text interpretation, an unsuitable statistical model, sampling error,
a missing variable or inconsistent restrictions supplied by different owners.
The proposed tool should return the smallest diagnosed conflicting subset when
it can compute one, and otherwise retain the full conflict. It should not
silently discard low-ranked edges until a DAG appears. The original evidence
and all relaxations remain visible. Producing a graph after changing the
restrictions is a new analysis whose identifying content has changed.

A soft restriction makes a different tradeoff. A score of the form
$S(G;D)+\lambda R_M(G;T)$ combines statistical fit with a model-derived
preference. The regularization weight $\lambda$ determines how much the
language preference can change the selected structure. Neither its numerical
scale nor an edge frequency from repeated prompts supplies a posterior
probability that the banking mechanism contains an edge. A meaningful
sensitivity analysis varies the strength of the preference, keeps the data
term and feasible graph family fixed, and traces which decisions change.
Soft-constraint discovery is a legitimate alternative to a hard restriction,
but it must be compared under its stated meaning
\citep{v12_la_2405_13551,v12_la_2407_15073}.

\subsection{An observational equivalence bound for an assistant}

A short project derivation explains a limit that a more fluent assistant cannot
remove without new information. Suppose two structural models $\mathcal M_0$
and $\mathcal M_1$ generate the same observed law for every available record and
the same admissible documents, while their target effects are $\theta_0$ and
$\theta_1$, with $\Delta=|\theta_1-\theta_0|>0$. Let the assistant, numerical
code and critic together be any randomized procedure using only those records
and documents. Its internal randomization has the same distribution in both
worlds and contains no world-specific oracle. Then its reported interval $C$
also has the same distribution in both worlds.

If the procedure claims coverage at least $1-\alpha$ under each model, the two
events $\{\theta_0\in C\}$ and $\{\theta_1\in C\}$ have probabilities at least
$1-\alpha$ under that common report law. The elementary union bound gives
\begin{equation}
 \Pr\{\theta_0,\theta_1\in C\}\geq 1-2\alpha,
 \qquad
 \Pr\{\operatorname{diam}(C)\geq\Delta\}\geq1-2\alpha.
 \label{exec14:language-information-bound}
\end{equation}
The diameter statement also holds for a set-valued report. It does not require
normality, a particular discovery algorithm or a claim that an LLM is
statistically calibrated. It says what uniform coverage would require in this
two-world experiment. For nominal 95 percent coverage and an effect separation
of $2/5$, diameter at least $2/5$ is necessary with probability at least $9/10$.
An algorithm that always returns a narrower interval cannot satisfy both
coverage claims. A refusal is an allowed scientific result; counting it as a
successful narrow interval would change the claim.

The same argument has a decision form. If the correct binary action differs
between the two worlds, let $q$ be the probability that the workflow selects
action one. With unit loss for a wrong action, its risks in the two worlds are
$q$ and $1-q$, in one order. Their maximum is at least $1/2$. More independent
prompt draws can estimate $q$ more precisely; they cannot make both risks small
when the observations and applicable text are the same. The way to improve the
decision is to obtain a discriminating restriction or intervention, accept a
robust action, or state the unresolved risk. The source of a useful restriction
may be a person or a document, but its validity is additional information.

This limitation is directly relevant to the funding equivalence example in the
Java adapter discussion, where several Gaussian models share an observed
covariance yet differ in the intervention response. It also sharpens a fair
language evaluation. If the supplied narrative explicitly tells the model which
world generated the data, the task assesses use of that information. It must
not be described as identification from the numerical observations alone. If
model pretraining contains relevant proprietary policy information, that
information is also part of the effective input and raises a provenance and
transport question; it is not an invisible extension of a statistical theorem.

\section{From edge accuracy to the bank's decision loss}

The discovery target must be defined before a leaderboard is useful. A DAG,
a CPDAG, a PAG and an RFCI-PAG make different assertions. A DAG adds directions
that may be unresolved in a CPDAG. A PAG can retain circles because ancestral
relations differ across compatible latent-variable models. An RFCI adjacency
can have a weaker meaning than an FCI adjacency. Comparing all four as lists of
directed edges against a single chosen DAG can reward unwarranted certainty
and punish scientifically correct ambiguity
\citep{exec13_latentr_colombo2012,exec13_latentr_primary}.

A proposed evaluation should therefore publish separate losses for the
skeleton, endpoint marks, invariant ancestral claims, extra assumptions and the
downstream estimand. If $E^\star$ is the true adjacency set under a declared
reference and $\widehat E$ the reported adjacency set, precision and recall
answer whether proposed links are right and whether reference links are found.
An orientation score computed only on $E^\star\cap\widehat E$ answers a
conditional question; it can look excellent when most links were missed.
Unknown endpoints should have their own accounting, with a penalty only for
claims the task actually requires. The interpretation of retrieval-assisted
orientation rates depends on precisely this denominator
\citep{v12_la_2402_15301}.

For an intervention, graph error matters through a loss function. Let
$\Theta(\widehat{\mathcal G})$ collect effects compatible with a reported graph
class and other stated restrictions. A proposed action can maximize worst-case
utility, maximize utility under an explicitly elicited model distribution, or
request more information. These are different decision rules. The assistant
must not silently replace the first with the second by selecting the graph it
finds most plausible. The nominal effect interval and the range over graph
models are also different objects; combining them requires a stated procedure.

Consider an exact planning example in which an offer costs 3 units per account
and its incremental gross return is either 0 or 8 under two observationally
compatible models. The net value of offering is either $-3$ or $5$, whereas
not offering has value zero. The worst-case expected-value rule selects no
offer. If an elicited probability $q$ is placed on the high-return model, the
expected net value is $8q-3$ and offering is chosen only for $q>3/8$. This
threshold makes the consequence of the extra belief visible. A language-model
vote frequency cannot be substituted for $q$ without a model relating votes to
which world is true. The costs and returns are stipulated units, not bank
prices or measured ROI.

An ideal, free observation that perfectly reveals the model would yield value
$5q$, since the offer would be made only in the high-return world. The value of
that information before paying for it is
\[
 5q-\max(0,8q-3).
\]
At $q=1/2$ it equals $3/2$ units per account. A proposed experiment costing more
than that amount per account could not be justified by this particular perfect
information calculation, and a noisy experiment cannot exceed its value under
the stipulated model. This is a planning bound, not a recommendation to conduct
a bank trial now. It identifies what an inexpensive public or generated
prototype should measure: whether the workflow preserves the ambiguity and
correctly reports the threshold for a decision.

The criterion changes when loss includes customer harm, regulatory constraints
or uneven effects across groups. Those obligations must enter the feasible
action set and utility definition explicitly. A large mean return cannot by
itself authorize an infeasible offer. Similarly, a loss expressed per account
requires care when one customer holds several accounts; the intervention and
interference units may differ. The graph workflow should surface these choices
before performing a cost comparison, rather than use the presence of a currency
symbol as evidence that the causal problem has been economically specified.

\section{Two channels of uncertainty and a calibrated refusal}

Graphical ambiguity and statistical uncertainty should survive into the same
report without being conflated. The first arises when the same population law
is compatible with more than one target effect under the stated information.
The second arises because that law is estimated from a finite sample. A third
source, uncertainty about whether the variables and assumptions are meaningful,
is not automatically represented by either a confidence interval or a PAG.
A language assistant is often most useful on this third source: it can find an
ambiguous definition or an unsupported restriction. That usefulness does not
make its confidence score a measure of the first two.

One proposed report therefore has a target object, an identification object,
a sampling object and an evidence-review object. The identification object
states ``identified under these assumptions'', ``nonidentified in this model
class'', or ``unknown because the available procedure or input is inadequate''.
Failure to find a back-door adjustment set with a restricted implementation is
not a proof of nonidentification. The sampling object states the estimator,
method of uncertainty calculation, population and any trimming or weighting.
The evidence object links assumptions to inspected sources and distinguishes
missing evidence from a retrieval failure. A failed numerical test is neither
independence nor evidence of an unobserved common cause.

CausalDS is a useful comparator because it makes abstention on specified
nonidentified tasks part of the evaluation, and separates conceptual variables
from their noisy observation layer. Its finite benchmark findings remain
reported evidence for those tasks. The source-level adaptations needed for our
prototype are explicit three-state identification metadata, finite output
checks and denominators that distinguish principled abstention from missing
files, parsing errors and numerical failures
\citep{v12_la_2607_08093,exec14_estimation_sources}.

A refusal can be correct without being useful enough for a product. The
comparison should score whether it names the missing assumption, unsupported
operation or conflicting clock and whether a feasible next action is supplied.
A rule that refuses every request avoids some false causal claims while failing
all identifiable tasks. A rule that returns a number for every request has high
completion while failing nonidentified ones. Plotting error and returned-answer
rate over a frozen threshold is more informative than a single completion
percentage, provided the threshold was selected separately from the held-out
assessment. All attempted requests, including refusals and malformed outputs,
remain in the operating-cost denominator.

A critic can select which analyses reach the report, but acceptance by a critic
is not automatically conditional statistical coverage. As an exact counting
example, six covered intervals among six accepted requests, with two refusals,
give selected coverage one and an all-request return-and-cover rate of $3/4$.
Neither denominator alone shows a nominal conditional guarantee. To investigate
that claim, a prospective experiment must define the selection rule, tasks and
repeated sampling law, and report uncertainty for its empirical rates. The
critic's access to a reference answer must be separated from the estimator's
information. Reviewing one's own answer with the same model adds an interaction,
not independent scientific evidence.

\section{Recent agent systems as competing scientific workflows}

Recent systems should be compared by the boundary at which the agent's output
becomes authoritative. In one architecture the agent directly proposes edges.
In another it chooses settings for a numerical engine. In a third it edits a
shared data model or proof statement. A fourth retrieves evidence for a human
reviewer. All can be described as causal agents, but they place different
obligations on the scientific core and on the person who accepts the result.
The banking prototype should choose a boundary that can be tested with the
information it actually has.

Zhu and colleagues' CausalAgent combines data processing, structure learning and
reporting with MCP tooling and a Sachs demonstration. The inspected paper places
quantitative effect estimation and counterfactual modules among its future
directions. Its strongest immediate contribution to this proposal is a
workflow pattern for moving from data and a question to a structured report.
The proposal cannot use the presence of MCP tools as evidence that an early
arrears-review effect has been identified, or that the future quantitative
module is already available. A direct integration decision must inspect the
called implementation and what it does when parsing or a backend call fails
\citep{v12_la_2602_11527,exec13_data_causalagent_code}.

Han and colleagues' separately authored Causal Agent is a different source.
It illustrates named graph objects stored in tool memory and a generated
question benchmark. This is useful for persistent state: an analyst should be
able to distinguish the graph before a restriction from the graph after it,
and the target before trimming from the target after it. Yet storing a graph
under a durable name does not validate its construction. The proposed service
binds a graph's content, variable definitions, information clock and added
assumptions together. Two graph objects with the same edge list can still be
different scientific inputs if their variables or populations differ
\citep{exec13_census_causal_agent}.

ORCA places more emphasis on relational state, including joins, variable roles,
graphs and estimation configuration. This addresses a real banking problem:
joining an account table to a customer table can multiply rows, change the unit
and create a misleading apparent sample size before discovery begins. Its
reported PAG-to-DAG step, however, makes additional orientation or editing
choices. Those choices must be visible as assumptions. The inspected report
also contains a discrepancy between the abstract's 42-percentage-point
completion improvement and a table with completion rates $.83$ and $.75$.
Completion is defined there by production of an ATE without requiring accuracy.
The proposal retains the discrepancy and avoids treating either a completion
figure or the existence of a final estimate as evidence of correct causal
inference. No corrected experiment is invented
\citep{exec13_census_orca}.

The causal-learn+ position puts assistance around a protected discovery core.
Agents can help inspect data, explain algorithms and interpret results while
formal algorithms retain responsibility for the numerical discovery operation.
That division is the preferred initial design here because it makes the
scientific method and the agent's added assumptions separately inspectable.
It still requires a correct core: a well-typed call to a flawed algorithm is
not a sound discovery result. It also leaves measurement, latent structure and
sampling assumptions open to review. The paper's call for future evaluation
supports a candidate architecture, not a measured banking advantage
\citep{v12_la_2606_23608}.

CausalSmith adds a different kind of protection by making a formal statement,
its proof dependencies and compiler outcome explicit. A checked proof answers
whether the formal statement follows within its logical environment. It does
not settle whether ``recovery'', ``offer'' and ``available before treatment''
were formalized faithfully, or whether the assumptions describe the bank.
The inspected limitations retain unvalidated model judgments about statement
faithfulness and research quality. The appropriate use is to expose a proof
obligation and its assumptions, then review the mapping to the business problem.
It is not to treat successful compilation as a substitute for that mapping
\citep{exec13_census_causalsmith}.

These choices lead to a concrete reuse recommendation. Adapt stateful
specification and evidence-reporting ideas, retain a fixed independently
reviewed numerical core, and require typed proposals for changes to the target,
variables or graph restrictions. Defer unrestricted autonomous execution until
the full package and operational path is audited. Reject automatic conversion
of an uncertain graph into an asserted DAG merely to enable a downstream
estimator. Keep proof assistance as a separate declared capability whose
formal assumptions remain inspectable. The proposal claims that these boundaries
make errors easier to observe; an empirical claim that they reduce review cost
or improve accuracy awaits a matched study.

This decision is revisable. A more autonomous system could be preferable if a
fixed public or generated evaluation shows that it preserves targets and graph
uncertainty with less total effort. A protected core could be ineffective if
most errors occur in the semantic mapping before the core is called. Both
outcomes are plausible hypotheses. The next experiment should therefore measure
where corrections occur and what information was needed to make them, rather
than count the number of agent roles or reward the longest explanation.

\section{Following the full package path rather than the package name}

A research package is not interchangeable with an executable operation of the
proposed library. The first path is a selected PC call in causal-learn, its
returned graph representation, a conversion into the shared graph schema,
a context-aware diagnostic report and a bounded CLI or MCP response. The
inspected MIT-licensed revision offers a plausible direct integration route,
but its twelve direct dependency declarations do not constitute a full pinned
environment. Porting the algorithm now would add maintenance responsibility
without a measured embedding constraint that justifies it. Executable admission
is therefore deferred while this path is completed; the reuse choice itself is
concrete and reviewable \citep{exec14_causallearn_reuse}.

The PC and FCI profiles must remain distinct. A PC bidirected conflict under an
orientation rule cannot silently become the latent-common-cause assertion of a
PAG. The conversion needs the two proximal endpoint marks, a bijection from
input variables to nodes, and isolated nodes. Unsupported composite endpoint
encodings remain unsupported. A Fisher-Z failure due to a singular inverse,
an invalid $n-|S|-3$ term or a nonfinite numerical value cannot become a valid
independence decision. An $\alpha$ threshold for one query does not supply
whole-graph or effect-set coverage. These are requirements for the first scoped
adapter, not benefits already delivered by listing causal-learn as a dependency
\citep{exec14_causallearn_reuse}.

Tetrad provides a materially different Java path. At the inspected revision,
FCI initializes circle endpoints and invokes orientation code with a flag for
excluding selection bias. The class-level reference to a CPDAG is inconsistent
with the PAG-oriented search path and cannot define the interchange contract.
RFCI is a separate procedure with its own graph meaning. The Edge constructor
can reverse the stored order of nodes for a left-pointing edge, so conversion
must bind each mark to its proximal node, not assume that construction order
is stable. Star and null endpoint values exist in the Java enum but are outside
the proposed initial PAG interchange profile
\citep{exec14_java_latent_sources}.

Static inspection also identifies a specific reason to withhold full FCI
behavioral qualification at that revision. In the inspected possible-D-SEP
block, the code builds one candidate set from each endpoint, with a path-length
argument of three, and tests that complete set. That block does not enumerate
all of its subsets. Conditional independence is not monotone in the conditioning
set: adding a collider can turn independence into dependence. The existing
exact Gaussian counterexample demonstrates this logical point; it is not an
executed counterexample to the entire Tetrad algorithm, whose earlier skeleton
phase and later rules can affect the actual candidate graph. The required next
check follows the whole search on an oracle fixture that reaches this block
\citep{exec13_java_examples,exec14_java_latent_sources}.

A second branch matters when configuration is changed. The ordinary
\texttt{finalOrientation} call is inside the possible-D-SEP branch in the
inspected search method. Setting that option false skips that call; a separate
optional guarantee-PAG branch is not a substitute for an unexamined claim of
algorithmic equivalence. The conclusion is bounded: the source path needs
qualification, not that every Tetrad release or every returned graph is wrong.
The available local Java 11 runtime and absent compiler also do not meet the
Java 21 build profile previously reviewed. No Java FCI, RFCI or endpoint
roundtrip has been executed in this review
\citep{exec14_java_latent_sources,exec13_tetrad_reuse_primary}.

The proper alternatives are an isolated, pinned Tetrad environment with
method-specific oracle and numerical tests; the separately reviewed R pcalg
FCI/RFCI routes; and a separately reviewed Python FCI profile. Their settings,
selection assumptions, graph classes and conversion conventions have to be
matched before any speed comparison. A benchmark that times a truncated search
against a fuller search without labeling the difference measures a change in
scientific work as well as an implementation difference. A process bridge adds
startup, serialization, memory and recovery costs that belong in the complete
workflow comparison.

\section{A hidden-risk example with identical observable distributions}

The information bound becomes tangible in a nondegenerate Gaussian example.
Here the banking action is a standardized concession intensity, rather than the
binary offer used in the earlier motivating brief. Consider two centered models.
In the first, independent disturbances have variances
$\operatorname{Var}(U)=1/2$ and
$\operatorname{Var}(\epsilon_A)=\operatorname{Var}(\epsilon_Y)=1/2$, and
\[
 A=U+\epsilon_A,\qquad Y=U+\epsilon_Y.
\]
In the second, $A$ has variance one and is independent of $\eta$, with
$\operatorname{Var}(\eta)=3/4$, and
\[
 Y=\tfrac12 A+\eta.
\]
Both produce the bivariate normal law with variances one and covariance $1/2$.
Both have a positive observed density on the plane. Their observed regression
of $Y$ on $A$ has slope $1/2$. In the first model, intervening on $A$ leaves
$Y$ unchanged, so the derivative of the intervention mean is zero. In the
second it is $1/2$. Every statistic of the observed pair has the same sampling
distribution in the two models for every sample size. Recovering a precise
association does not decide which intervention response is right.

The example isolates the role of private borrower information. The first model
contains a common cause that is unmeasured in the supplied table; the second
attributes the association to the concession. An agent given only the same
table and an ambiguous description has no model-specific evidence that selects
one. Saying that the concession ``obviously affects recovery'' introduces an
assumption. It may be reasonable as part of a scenario, but it must be exposed
as such. A PAG over two observed variables can retain unresolved endpoint marks;
returning a single oriented edge removes an uncertainty that the data did not
remove.

Suppose a new predecision score $W=U+\epsilon_W$ is offered as a remedy, with
independent noise variance $s$. In the common-cause model the residual variance
of $U$ given $W$ is
\[
 v=\frac{(1/2)s}{1/2+s}=\frac{s}{1+2s}.
\]
After residualizing $A$ and $Y$ on $W$, their covariance is $v$ and the residual
variance of $A$ is $1/2+v$. The adjusted coefficient is therefore
\[
 \beta_{A\mid W}=\frac{v}{1/2+v}=\frac{2s}{1+4s}.
\]
For $s=1/2$ it is $1/3$ although the true causal response is zero. With an exact
score, $s=0$, it becomes zero; as the score becomes uninformative it approaches
$1/2$. This is a separate model from the unit-variance text-measurement example
in the earlier chapter. The variance choices and action definition are stated
so that the two examples are not mistakenly pooled. A document extractor can
help create $W$, but the identifying question is what information its error
model and validation sample supply.

The remedy is specific. Obtain a defensible measurement model with information
sufficient for the proposed correction, a justified assignment restriction, or
an intervention that distinguishes the worlds. A smaller standard error from
more observations of $(A,Y,W)$ does not remove the residual coefficient when
$s$ stays positive. A language agent can explain the alternatives, request the
necessary metadata and prevent a false report of identification. It cannot
invent an independent validation sample by generating plausible labels.
The exact covariance and residual-variance calculations in this section are
project derivations, reproduced in rational arithmetic
\citep{exec14_discovery_oracles}.

\section{Why a CI threshold is not graph or effect coverage}

A discovery interface often exposes one number called alpha. Its scope must be
made explicit to an agent and to a business reviewer. For a particular fixed
conditional-independence test, a level-$\alpha$ statement concerns the
probability of rejecting that null when it is true under the declared sampling
law. It does not bound the probability that a whole graph is correct, that all
orientations are right, or that a derived effect interval covers its target.
A graph can fail because a false independence is accepted, an assumption is
wrong, a query is numerically invalid, or an orientation is misinterpreted.
Those errors are not all controlled by a single null-rejection level.

Even the limited false-rejection claim changes when many queries are possible.
For $p$ observed variables, all unordered variable pairs and all conditioning
subsets of the remaining variables form a family of size
\[
 M=\binom p2 2^{p-2}.
\]
If each fixed query has unconditional type-I error at most $\alpha$, then a
union bound controls at least one false rejection among that fixed finite
family by $M\alpha$, truncated at one. No independence between query decisions
is needed for this inequality. It applies to any adaptively selected subset of
the family only when the procedure uses the same prespecified tests and the
per-query calibration holds for every fixed query; an adaptive prompt that
changes the data or test definition is not covered by that argument.

With six variables there are 240 possible pair/subset queries. At $\alpha=.05$
the bound is uninformative. At $\alpha=.0002$ it is $.048$ for the specified
false-rejection event. This second number is still not a 95.2 percent guarantee
for the graph: it does not control missed alternatives, faithfulness violations,
latent-variable misspecification or endpoint conversion. The example is useful
because it gives the assistant a precise explanation of what alpha can and
cannot mean. It is not a recommendation to choose a universal threshold of
$.0002$ for bank discovery.

The same distinction applies to language-response tests. Their sampling law is
a law of model answers under a prompt and decoding configuration. A response
frequency can be estimated under assumptions about repeatability, but it does
not inherit the error guarantee of a numerical CI test. An agent that reads a
column named ``pvalue'' must check the generating procedure. The response-based
variance calculation in the preceding material demonstrates one implementation
hazard; the more fundamental issue is that the response null and the banking
distributional null are different statements.

For the prototype, every CI result should therefore include the query, sampling
unit, conditioning variables, test method, usable observations and computation
status. A singular matrix produces an unavailable test, not a very strong
scientific dependence claim. The search result should state whether all
required queries completed and whether any search bounds or time limits
changed the explored family. The report should retain sensitivity to thresholds
and to permitted restrictions without presenting that sensitivity exercise as
a calibrated confidence set. A qualified effect procedure needs its own
identification and uncertainty argument after the graph step.

\section{A public evaluation design that can change the proposal}

The immediate evaluation should test whether language assistance improves the
causal workflow within a fixed information budget. It can use public documents
for realistic vocabulary and project-generated mechanisms for known truth,
with the two origins clearly separated. Generation does not create authentic
bank assignment evidence, but it makes intervention effects, measurement error,
latent variables and timing failures observable to a held-out scorer. Private
truth must not enter the assistant's prompt, retrieval corpus, tool output or
file tree. The predeclared deterministic baseline receives the same permitted
brief and data as the assisted workflow.

A useful task portfolio crosses four dimensions. One varies whether the effect
is identified from the supplied graph and observations. Another varies whether
variable names are informative, neutral or misleading. A third changes whether
the document's reported time and its availability time agree. A fourth varies
whether text contains a valid additional restriction, a contradiction or no
relevant information. This is a design proposal; choosing the generator, sample
size, splits and replication count requires a prospective protocol before any
performance study. A successful run on a few authored examples is a conformance
check, not a retrospective preregistration.

The banking tasks should include an offer distinct from uptake, a branch-level
intervention with account-level observations, an unrecorded borrower-risk
variable, and a postdecision note that looks like a baseline covariate. They
should include a graph with correct but unresolved endpoints and a case where
an apparent numerical failure is caused by exact accounting identities. A
language model that recommends a valid operation but uses the wrong treatment
variable has not answered the task. Conversely, a workflow that retains the
right target and returns a justified nonidentification result can succeed on a
case for which no point estimate is licensed.

Names and mechanisms should be perturbed separately. Replacing all banking
terms with anonymous labels can remove legitimate intervention and timing
information, so that perturbation is scored as a different information regime.
Holding the definitions fixed while changing irrelevant wording probes another
kind of robustness. Reversing a generated mechanism under familiar names tests
whether the model follows data or a memorized story. The inspected work on
pretraining relation frequency motivates this distinction, but its findings
for accessible models and commonsense relationships are not a theorem about
all language models or all banking questions
\citep{v12_la_2407_19638}.

Document-level leakage also needs a precise split. If the same policy is
paraphrased into both development and evaluation sets, a row-level random split
does not produce an unseen policy. If several accounts share one assignment
rule or one customer, independent-row formulas may overstate the effective
sample. The proposed protocol should group related documents, mechanisms and
customers according to the claim it seeks to assess, and say which kind of
transfer remains untested. No generic instruction to ``hold out 20 percent''
settles these scientific dependencies.

For each request the result should contain the target identity, permitted
information, proposed graph class, assumptions added by the agent, formal
identification state, numerical method and final disposition. Metrics should
include false causal assertions, missed identifiable opportunities, unsupported
orientations, effect error where meaningful, interval loss and all-request
resource cost. Model-only latency omits retrieval, serialization, review,
retries, corrections and refused jobs. A cheap request that repeatedly needs an
analyst to repair a target can cost more than a slower fixed workflow.

The proposal changes if the assisted path reliably reduces a specified review
burden while maintaining these scientific boundaries on a frozen task family.
It also changes if assistance adds unsupported restrictions or hides failures,
even when the prose improves. A later live comparison must bind provider and
model versions, prompts, decoding, tool identities, custody and actual prices;
none of these results is inferred from this design. An expensive bank experiment
becomes relevant only after the prototype isolates an uncertainty that public
or generated evidence cannot answer and shows why resolving it can alter the
investment decision.
\label{exec14:l4-development-end}
